\chapter[Sermon 47. Of the Two Prophets Fighting Manfully Against Antichrist, and of Their Power.]{Of the Two Prophets Fighting Manfully Against Antichrist, and of Their Power.}
\label{sermon:047}
\markright{Sermon 47. Of the Two Prophets Fighting Manfully Against ...}

And I will give power to my two witnesses, and they shall prophesy one thousand two hundred sixty days clothed in sackcloth. These are two olive trees, and two candelabra standing before the God of the Earth. And if any man will hurt them, fire shall come forth from their mouth and devour their enemies. And if any man will hurt them, this way must he be killed; these have power to shut heaven, that it reign not in the days of their prophesying; and have power over waters to turn them to blood, and to smite the earth with all manner of plagues as often as they will.

\index{Antichrist}These things also pertain to the comfort of the faithful. For the Lord promises that he will send prophets—that is, preachers—who will maintain and defend the truth of the Gospel and the glory of Christ, confront the Antichrist, and destroy his kingdom, and advance the salvation of the faithful. In the previous chapter, verses 8 and 9, the conflict of the Antichrist and heretics against God and his Christ, and against his church, was described; and now, in a few words, the opposing conflict is set forth, and the army of Christ is mustered.

\index{Jerome, Saint}And he brings forth two prophets, that is, preachers: not that there shall be only two, but to signify that the power of Christ in the world would appear small to worldly men (as I shall tell you anon). Meanwhile, he understands all faithful preachers and pastors of all times who offer themselves to resist Antichrist and heretics. Some interpret these things as referring to Enoch and Elijah, who will appear physically before the judgment. However, Saint Jerome in his letter to Marcella refers that opinion to Jewish fables, signifying that these things must be understood spiritually, as are most things in this book. And, generally, all interpreters agree in interpreting these things as spiritual prophets, not physical ones according to the literal sense. I suppose there are only two prophets mentioned here for two reasons.

\index{Daniel (Book of)}\index{Zechariah (Book of)}First, for that he would allude to the old history or prophecy of Zechariah, which is in chapter 4. It was thought then, also, that the people of Israel, returned from Babylon, would find the rebuilding of the Temple impossible, for they had many and mighty adversaries, and they were weak and few, and their governors, Zerubbabel and Joshua, were contemned. But through the mighty hand of God and his faithful aid, it came to pass that the power of their adversaries vanished away as vain, and they, in defiance of hell gates, built up their Temple right, so the Lord says it shall be in that later age, that the ministers, most contemptuous and very few in number, shall build up Christ his temple and repair it, and shake the most mighty power of Antichrist. Hereunto I suppose belongs that saying of Daniel: and when they shall fall, they were helped with small aid, \&c. Secondly, for this cause chiefly he accounts only for two witnesses, for it is written in the Law, “in the mouth of two or three witnesses every word shall stand.” It is judged therefore a full testimony, which shall be confirmed with the agreeable declaration of two. Wherefore, when the Lord says that he will give two prophets, it is as much to say as that he will give so many ministers as shall suffice, which shall both build up his church and also pluck down and restrain the kingdom of Antichrist. Some of the commentators think that by two witnesses are understood two testaments. However, we see that the Lord speaks here of witnesses, not of the thing testified or to be witnessed, which nevertheless we do not separate from the witnesses.

\index{Papacy}The apostles and those considered apostolic are called witnesses everywhere in the Gospel and in the first chapter of the Acts of the Apostles. Witnesses are appointed in judgment to faithfully declare what they have seen or heard, to invent nothing of their own, and to neither add nor subtract anything from what is testified. Similarly, God places witnesses of God—that is, ministers—in the church of God, and it is required of them that they imagine nothing of their own, nor add to or take away anything from God’s word, but simply declare to the church of God the things they have seen in the narrative of the Gospel and heard from the prophets and apostles. Therefore, those who do not bring the Gospel are false witnesses, unworthy to be called witnesses of God and of Christ; they are rather the Popes’ witnesses, whose decrees and decretals they bring forth and testify to for the unlearned people. Therefore, those two prophets will be witnesses of Christ and will bring witness for Christ from the most true Scriptures.

And the beginning of these things is here referred to God and to his Christ, as the origin of Antichrist is reduced to the devil himself. “I will give,” says the Lord, “to my two witnesses, and they shall prophesy.” Christ sends preachers, and gives to them also the ability to preach. This is a wonderful comfort. For just as the devil often sends, instructs, and helps his false prophets, so Christ does not leave his church destitute, but gives to these ministers the ability to teach and do things successfully. For in the Gospel also he promised and said: “I will give you a mouth and wisdom, which they shall not resist, as many as are against you.” These things ought to comfort us in the grievous consultations, betrayals, and assaults of the enemies of the Gospel. Christ will not forsake his ministers, so long as they are faithful and depend upon Christ alone.

Now is also declared the time of the preaching of the gospel against Antichrist, truly all that time wherein Antichrist shall tread the Temple and holy city. For a thousand two hundred and sixty days make forty-two months, if you put thirty days to every month. But we heard before that Antichrist should tread the church forty-two months. Again therefore, a certain number is put for an uncertain [?]. And here is signified, and that with a mystery is here defined the time of days, not of months or years. For though the function of the ministry be never so hard and dangerous, yet so shall God comfort and confirm them, that they may appear a few days only, not months or years to suffer persecution, and to travel in this laborious work of the Lord. And where I have said that those numbered days are put for an uncertainty of time, this has moved me, that by and by in the twelfth chapter, the same number of days shall be assigned; for which yet he has set before, for a time, and times, and half a time. Which appears plainly to be taken out of the seventh and twelfth chapters of Daniel. I know that the same is expounded by many for three years and a half; that the time should signify a year, times, two years, and half a time, half a year. But every man may perceive that the thing itself is repugnant to that number of years, if he be at the least any thing seen in histories. In the seventh chapter of Daniel, he says, the other beasts gave over their rule, and spaces of life were granted, for a time, and a time. But who will expound these things of two years only, since it is evident that the Babylonians, Persians, and Macedonians reigned many years? He signifies therefore that those kingdoms should reign so long as God would permit them, and give them power to reign. We say in Dutch, where yet we appoint no time prefixed. In the same chapter of Daniel is put the same phrase of speech, that the saints shall be delivered into the hand of Antichrist, for a time, times, and half a time. And in the twelfth chapter, he says that his prophecy shall be fulfilled in a time, times, and half a time. But who shall believe that within three years and an half all those things should be accomplished, which he declared in the whole work? Why then do they restrain the times of Antichrist to three years and an half, especially his persecution? Why see they not the destruction of Antichrist, and the peace of saints, and the day of judgment, to be the same day? For Daniel says, that the beast should be cast down headlong into Hell, when the seats be furnished. And Paul says, whom he shall destroy with his coming, and who shall show unto us the certain day of judgment? It is known to the Father alone. Let them leave therefore with their computations to strive with the Gospel. It appears therefore that the Lord by that kind of speaking, as it were by a riddle, has defined no time certain; but rather has admonished the godly of long suffering, of patience, and constancy; and has commanded that we should not over curiously search into the instant of this time, but should rather permit it to Christ himself, in another place saying: It belongs not to you to know times, and the moments of times, which the Father has reserved in his own power; but watch, that when the Lord shall come, he may find you watching. Therefore whether so ever the Lord shall defer his judgment a long, short, or mean time, be you constant. So at this present he says, how the ministers of Christ shall preach all that time, wherein Antichrist shall persecute. And truly, if thou read the histories, thou shalt find that the most virtuous and best learned men, have in all ages, now for the space of these seven hundred years and more, constantly resisted the Popes enterprises, their great abominations, and crafty jugglings and seductions of monks and Friars. Of the persecutions that they have suffered, I will speak hereafter.

\index{John, Saint}\index{Repentance}Furthermore, the appearance of these prophets is also shown, so that the manner of their doctrine may be understood from it. They will not be clothed in soft or precious garments, such as velvet, satin, or damask, or crimson embroidered, but in sackcloth. And sackcloth, as appears in the Prophets, is for a mourning garment and for those who are penitent. Therefore, just as Saint John was plainly clothed and preached repentance, so these also will urge repentance and amendment of life, and persuade people to frugality, while persecuting riot and all intemperance. Certainly, all good and learned men for these past seven hundred years have desired nothing else of the Pope and clergy, and of the people, but repentance and reformation; for which they have received little thanks from them. But what the apparel of the Antichristians is, no one is ignorant at this day. Some of it does not differ much from that of harlots. Consequently, he declares more fully and at length, of what sort they will be, and also their ministry, what the effect and virtue of their preaching will be. And he sets this forth and declares it with various figures taken from the scriptures.

First, he alludes again to the fourth chapter of Zechariah. These are two olive trees, and lamps are nourished with oil; therefore, oil signifies the matter of preaching or of sermons. For candelabras bearing lights are preachers, showing abroad the light of Christ and his gospel throughout the world. And that preaching of light is taken from Scripture, just as the light of a candelabra is nourished with oil. Oil is a sign of the holy of holies. Therefore, John also calls the Holy Spirit “unction.” Certainly, Scripture is the inspiration of the Holy Spirit. Therefore, those preachers shall preach Christ from Scripture. And so, preaching the gospel of Christ through the inspiration of the Holy Spirit, they are said to stand before the sight of God of the earth: that is to say, they are in the protection, the care, and providence of that God by whose providence are governed whatsoever are in heaven or on earth.

For he appears to have alluded to these words of Zechariah: “The eyes of the Lord look over the whole earth,” and these are the two children of oil, which stand before the governor of the whole earth. And these things exceedingly comfort faithful preachers, who see that God has care of them, I mean God the Lord of all.

Again, they are neither olive trees nor lampstands displaying the light of the Gospel, but those of the Antichrist’s party consider them the dregs and filth of men, placing them in the stead of the oil of the Holy Spirit and putting them into the lamp. Neither do they display any light, but darkness and the opinions of the most corrupt men. Against these, John reasons, saying, “These things I have written to you concerning those who deceive you.” And the unction which you have received from him remains in you, and you have no need that anyone teach you; but just as the very unction teaches you all things, so is it true, and no falsehood.

Now also the weapons of these preachers are described, wherewith they may defend their cause and fight against their enemies. If any man will hurt them, fire issues out of their mouth and devours their enemies. This signifies, with pretended malice and against justice to hurt or injure: and first he said to hurt. If therefore any of the champions of Antichrist shall assail those preachers and shall blame their doctrine and ministry, straightways they shall bring forth of the holy scriptures God’s word, and so shall repress and overcome their enemies. For that these things may not be expounded after the letter, that same chiefly proves that by and by we shall hear that those prophets shall be vanquished and put to death by Antichrist: to wit, corporally. Who then can not gather hereof that the victory of preachers is spiritual, that their adversaries vanquished of the verity may live in deed bodily, but through the virtue of the verity they may seem to be ghostly slain? And therefore as it were by an interpretation is added: and if any will injure them, so must he be slain. So I say, by fire verily which goes out of their mouth. And who will say that material and natural fire should come forth of a man’s mouth? And Paul also expounding these things, taking the manner of speaking of Isaiah, reasoning of Christ and of Antichrist: whom he shall kill, says he, with the breath of his mouth. Behold, Paul calls it the breath of the mouth, which John named fire. We read also in the twenty-third chapter of Jeremiah, “Is not my word as fire, and as a mallet breaking the rock?” And again in the fifth chapter, “As much as you speak this word, behold I will make my words in your mouth fire, and this people wood, and it shall consume them.” Of Elias we read in the fourth chapter of Kings, first chapter, that calling down fire from heaven he bodily burnt the king’s servants. Which example where the disciples James and John alleged, the Lord forbad them, that he might admonish them of their function, to wit, that they must fight with long suffering and with the word of the verity. Which the Apostle in another place commands expressly, to wit, in Second Timothy, second chapter. Whereby we are plainly taught that Antichrist must not be vanquished with corporal weapons by the ministers, but with spiritual. For he must be slain with the gospel, with that most sharp sword, and fall down and die in the breasts of men, that he may be utterly contemned and known to be Antichrist. And where many confound the ministry of the word and the power of the magistrate, and for the same cause take the sword out of his hands, commanding that in this case he may not strike heretics and blasphemers, affirming that they ought not otherwise to be punished than by the word: let them learn to discern better betwixt offices, and not to give that liberty to blasphemers and to all manner of seducers, and to such as having been a thousand times convicted of heresy, cease not to infect innumerable and bring them into perdition, unless they be straightly punished by the magistrate. Let every one therefore apply their own office, and herein follow the rule of verity and equity, and then shall things be in better order.

Furthermore, he adds more explicit statements concerning their power and ministry, even alluding here to various passages of Scripture. First, he says they have power to shut heaven so that it does not rain during their prophesying. He alludes to the story of Elijah, which is read in 1 Kings 17. And these things must be spiritually applied to our situation. For just as Elijah, through the power of God, prohibited rain, so the preachers of the gospel will, toward the disobedient, or those who will not hear the word but prefer to be seduced by papal abominations, shut up heaven itself, that is, assuredly testify that it is shut off from God, for as much as through Christ alone, as the only gate by which it is opened, they are restrained; and they will also tell them sharply that the grace of God is denied them, which is only granted by Christ. For the prophets are authorities that rain signifies the grace of God and the fruitful watering sent down from heaven. Therefore, during the time of their prophecy, they will steadfastly testify that those who prefer the Pope’s doctrines to the true bread from heaven are, through their own fault, deprived of that heavenly grace, light, and life. And again, we understand that they have power given to them to open heaven to believers. There is no place here to speak of that now. The things written in the gospel concerning the keys of the kingdom of heaven are more manifest, and chiefly belong to this, than that I should now rehearse them; I have both at other times and before in this same book spoken of them at length.

Secondly, he alludes to the story of Moses, and says that power is given to these prophets to turn waters into blood, which contradicts nothing with the former point. For the water of godly wisdom is a figure of the grace and relief of the Spirit. Blood signifies offense and punishment. For that saying of the law and of the Apostle is well known: “Your blood be upon your own head.” Therefore, these prophets will testify that God has truly sent his word of salvation to save all believers, but that it will be to the unbelievers, through their own fault, unto condemnation. For those who hear the preaching of God’s word and do not believe it, hear it to their own condemnation. And so the gospel is preached at this day to many without fruit, being corrupted by the popish doctrine, and by force will not be wise, etc.

Finally, they have power to strike the earth with every plague, as often as they will. But they will not, except God’s word, by which they, being inspired and instructed, are governed, shall command them. For they will do nothing wilfully; they will not follow their affections, but the word of God. However, they are said to strike the earth with plagues when, out of God’s word, they threaten that God with plagues will punish the sins of men. Those plagues are recited in chapters 26 and 28 of Deuteronomy. Therefore, if they threaten impenitent persons with war, pestilence, famine, sicknesses, and other calamities, God will send them to those who are incurable, as the Lord himself says often in Jeremiah. Again, and on the contrary, they shall enrich with all blessings those who obey God’s word, when they show forth the Lord’s blessing.

Thus much he has spoken hitherto concerning the preachers of the Gospel, who shall fight against Antichrist in that last age before the judgment, and shall build up the church and confirm the believers. You yourself shall observe in what preachers you perceive these marks, and you shall acknowledge them as the lawful prophets of God. And you shall acknowledge with all how great a benefit of God it is to have true and faithful preachers of God’s word. The Lord our God confirm all ministers of his word in the setting forth of his truth, to the world’s end.
